% Part: first-order-logic
% Chapter: introduction
% Section: satisfaction

\documentclass[../../../include/open-logic-section]{subfiles}

\begin{document}

\olfileid{fol}{int}{sat}

\olsection{ارضا}

همین که بدانیم !!{formula}s چیستند، می‌توانیم از پیش به معناشناسی منطق
مرتبهٔ اول بپردازیم: تعریف پایه در اینجا تعریف !!a{structure} است. برای
زبان سادهٔ ما، !!a{structure}~$\Struct M$ فقط سه مؤلفه دارد: مجموعه‌ای
ناتهی چون $\Domain{M}$ که \emph{!!{domain}} نامیده می‌شود، چیزی که
$\Obj a$ در~$\Struct M$ مشخص می‌کند، و چیزهایی که $\Obj P$ در
$\Struct M$ درباره‌شان صادق است. شیئی را که~$\Obj a$ مشخص می‌کند با
~$\Assign{\Obj a}{M}$ و مجموعهٔ چیزهایی را که $\Obj P$ درباره‌شان صادق
است با~$\Assign{\Obj P}{M}$ نشان می‌دهیم. !!^a{structure}~$\Struct{M}$
فقط از همین سه چیز تشکیل می‌شود: $\Domain{M}$،
$\Assign{\Obj a}{M} \in \Domain{M}$ و
$\Assign{\Obj P}{M} \subseteq \Domain{M}$. حالت کلی پیچیده‌تر خواهد
بود، زیرا !!{predicate}s و !!{constant}s فراوانی وجود خواهند داشت،
!!{predicate}s می‌توانند بیش از یک موضع داشته باشند، و !!{function}s
نیز در میان خواهند بود.

این مقدار برای تعریف ارضای !!{formula}sی که شامل !!{variable} نیستند
کافی است. ایده آن است که تعریفی استقرایی و متناظر با شیوهٔ تعریف
!!{formula}s به دست دهیم. مشخص می‌کنیم یک فرمول اتمی چه هنگام در
~$\Struct{M}$ ارضا می‌شود، و سپس، برای نمونه، بر پایهٔ اینکه $!A$ در
~$\Struct{M}$ ارضا می‌شود یا نه، تعیین می‌کنیم $\lnot !A$ چه هنگام
در~$\Struct{M}$ ارضا می‌شود. برای نمونه، می‌توانیم چنین تعریف کنیم:
\begin{enumerate}
  \item $\Atom{\Obj P}{\Obj a}$ در~$\Struct{M}$ ارضا می‌شود هرگاه و
  تنها هرگاه $\Assign{\Obj a}{M} \in \Assign{\Obj P}{M}$.
  \item $\lnot !A$ در~$\Struct{M}$ ارضا می‌شود هرگاه و تنها هرگاه
  $!A$ در~$\Struct{M}$ ارضا نشود.
  \item $(!A \land !B)$ در~$\Struct{M}$ ارضا می‌شود هرگاه و تنها
  هرگاه $!A$ در~$\Struct{M}$ و $!B$ نیز در~$\Struct{M}$ ارضا شوند.
\end{enumerate}
فرض کنید $\Domain{M} = \{0, 1, 2\}$، $\Assign{\Obj a}{M} = 1$ و
$\Assign{\Obj P}{M} = \{1, 2\}$. این تعریف به ما می‌گوید
$\Atom{\Obj P}{\Obj a}$ در~$\Struct{M}$ ارضا می‌شود (زیرا
$\Assign{\Obj a}{M} =  1 \in \{1,2\} = \Assign{\Obj P}{M}$). همچنین
می‌گوید $\lnot \Atom{\Obj P}{\Obj a}$ در~$\Struct{M}$ ارضا نمی‌شود؛
در نتیجه $\lnot\lnot \Atom{\Obj P}{\Obj a}$ ارضا می‌شود و
$(\lnot \Atom{\Obj P}{\Obj a} \land \Atom{\Obj P}{\Obj a})$ ارضا
نمی‌شود، و همین‌طور ادامه می‌دهیم.

مشکل هنگامی پدید می‌آید که می‌خواهیم سورها را تعریف کنیم: دوست داریم
چیزی مانند این بگوییم که «$\lexists[\Obj v_0][\Atom{\Obj P}{\Obj
v_0}]$ ارضا می‌شود هرگاه و تنها هرگاه $\Atom{\Obj P}{\Obj v_0}$ ارضا
شود.» اما !!{structure}~$\Struct{M}$ نمی‌گوید با !!{variable}s چه
کنیم. آنچه واقعاً می‌خواهیم بگوییم این است که
$\Atom{\Obj P}{\Obj v_0}$ \emph{برای یک مقدار~$\Obj v_0$} ارضا می‌شود.
برای دقیق‌کردن این مطلب به راهی نیاز داریم که نه‌فقط به $\Obj a$، بلکه
به~$\Obj v_0$ نیز !!{element}sی از~$\Domain{M}$ نسبت دهد. بدین منظور،
\emph{گمارش‌های} !!{variable} را معرفی می‌کنیم. !!^a{variable} گمارش
صرفاً تابعی چون~$s$ است که !!{variable}s را به !!{element}s
از~$\Domain{M}$ می‌نگارد (در مثال ما، به یکی از $0$، $1$ یا~$2$).
چون از پیش نمی‌دانیم کدام !!{variable}s ممکن است در !!a{formula} ظاهر
شوند، نمی‌توانیم متغیرهایی را که $s$ به آن‌ها مقدار می‌دهد محدود کنیم.
راه‌حل ساده این است که بخواهیم $s$ به \emph{همهٔ}
!!{variable}s~$\Obj v_0$، $\Obj v_1$، \dots\@ مقدار دهد. ما فقط
مقادیر لازم را به کار خواهیم برد.

به‌جای آنکه ارضای !!{formula}s را فقط نسبت به !!a{structure} تعریف
کنیم، آن را نسبت به !!a{structure}~$\Struct{M}$ \emph{و}
!!a{variable} گمارش~$s$ تعریف می‌کنیم و به‌اختصار
$\Sat{M}{!A}[s]$ می‌نویسیم. تعریف ما اکنون شامل بندی اضافی برای
!!{formula}s اتمیِ دارای !!{variable} خواهد بود:
\begin{enumerate}
  \item $\Sat{M}{\Atom{\Obj P}{\Obj a}}[s]$ هرگاه و تنها هرگاه
  $\Assign{\Obj a}{M} \in \Assign{\Obj P}{M}$.
  \item $\Sat{M}{\Atom{\Obj P}{\Obj v_i}}[s]$ هرگاه و تنها هرگاه
  $s(\Obj v_i) \in \Assign{\Obj P}{M}$.
  \item $\Sat{M}{\lnot !A}[s]$ هرگاه و تنها هرگاه نه
  $\Sat{M}{!A}[s]$.
  \item $\Sat{M}{(!A \land !B)}[s]$ هرگاه و تنها هرگاه
  $\Sat{M}{!A}[s]$ و $\Sat{M}{!B}[s]$.
\end{enumerate}
بسیار خوب، این کار یک مشکل را حل می‌کند: اکنون می‌توانیم بگوییم
$\Struct{M}$ چه هنگام $\Atom{\Obj P}{\Obj v_0}$ را برای مقدار
~$s(\Obj v_0)$ ارضا می‌کند. برای آنکه تعریف
$\lexists[\Obj v_0][\Atom{\Obj P}{\Obj v_0}]$ درست باشد باید یک کار
دیگر انجام دهیم: می‌خواهیم
$\Sat{M}{\lexists[\Obj v_0][\Atom{\Obj P}{\Obj v_0}]}[s]$ هرگاه و تنها
هرگاه $\Sat{M}{\Atom{\Obj P}{\Obj v_0}}[s']$ برای \emph{یک} شیوهٔ
~$s'$ از نسبت‌دادن مقدار به~$\Obj v_0$ برقرار باشد. اما مقداری که به~$\Obj v_0$
نسبت داده می‌شود لزوماً همان مقداری نیست که $s(\Obj v_0)$ مشخص
می‌کند. نمادی برای این وضع معرفی می‌کنیم: اگر $m \in \Domain{M}$،
آنگاه $\Subst{s}{m}{\Obj v_0}$ را گمارشی می‌گیریم که (برای همهٔ
!!{variable}s جز~$\Obj v_0$) درست مانند~$s$ است، با این تفاوت که
به~$\Obj v_0$ مقدار~$m$ را نسبت می‌دهد. اکنون تعریف ما می‌تواند چنین
باشد:
\begin{enumerate}\setcounter{enumi}{4}
  \item $\Sat{M}{\lexists[\Obj v_i][!A]}[s]$ هرگاه و تنها هرگاه برای
  یک~$m \in \Domain{M}$، $\Sat{M}{!A}[\Subst{s}{m}{\Obj v_i}]$.
\end{enumerate}
آیا درست کار می‌کند؟ فرض کنید برای هر~$i \in \Nat$ بگذاریم
$s(\Obj v_i) = 0$. داریم
$\Sat{M}{\lexists[\Obj v_0][\Atom{\Obj P}{\Obj v_0}]}[s]$ هرگاه و تنها
هرگاه $m \in \Domain{M}$ای وجود داشته باشد که
$\Sat{M}{\Atom{\Obj P}{\Obj v_0}}[\Subst{s}{m}{\Obj v_0}]$. چنین
مقداری وجود دارد: می‌توانیم $m = 1$ یا $m = 2$ را برگزینیم. توجه کنید
که این مطلب حتی هنگامی درست است که مقدار~$s(\Obj v_0)$ که خود $s$
به~$\Obj v_0$ نسبت داده است---در اینجا، $0$---کارساز نباشد. داریم
$\Sat{M}{\Atom{\Obj P}{\Obj v_0}}[\Subst{s}{1}{\Obj v_0}]$، اما
$\Sat{M}{\Atom{\Obj P}{\Obj v_0}}[s]$ برقرار نیست.

اگر این تعریف گیج‌کننده و دست‌وپاگیر به نظر می‌رسد، واقعاً چنین است.
اما این پیچیدگی افزوده برای ارائهٔ تعریفی دقیق و استقرایی از ارضای همهٔ
!!{formula}s لازم است و برای تعریف دقیق مفاهیم معنایی به چیزی از همین
دست نیاز داریم. راه‌های دیگری هم هست، اما همگی به یک اندازه
(نا)ظریف‌اند.

\end{document}
